\domainevide{Les cotes de ce thème supposent le groupe fondamental profini des courbes et l'action de Galois sur les dessins d'enfants, absents de mathlib.}{81, 138, 140-1, 144, 146}
\noindent Pistes de l'onglet Findings pour ces cotes, hors de portée de Lean pour l'instant :
\begin{itemize}\small\setlength{\itemsep}{0pt}
\item Cote~138 : une \q{opération} sur les cartes orientées donnée par une partie $A$ de la sphère de référence est, sauf un cas exceptionnel, la composition avec un revêtement $\varphi$ ramifié au-dessus de $0$, $1$, $\infty$, de ramification $2$ au-dessus de $1$. Mathlib n'a pas les revêtements de Belyi de la sphère.\piste{138-operations-belyi-extending}
\item Cote~138 : critère pour qu'une partie $A$ de la sphère définisse par image réciproque une nouvelle structure de carte : $A$ connexe, et $A$ contient au plus un des points $0$, $1$, $\infty$ ; telle qu'elle est lue, la seconde condition est contredite par l'exemple $A = [0, 1]$. Mathlib n'a pas les cartes sur les surfaces ; la question relève d'abord de la lecture.\piste{138-operations-subset-criterion}
\item Cote~140-1 : pour $S$ régulier intègre avec $\mathrm{Pic}(S) = 0$, le groupe de Picard du champ quotient $[(\mathbf P^1 - \{0, 1, \infty\})_S/\mathfrak S_3]$ est $\mathbf Z/6$ en caractéristique générique $\neq 2$, et $\mathbf Z/3$ en caractéristique générique $2$. Mathlib n'a pas les groupes de Picard des champs quotients.\piste{140-1-equivariant-pic-u03}
\item Cote~146 : table des plus petits exposants qui rendent triviaux, de façon compatible au groupe symétrique, les fibrés $L_i$ et $L = \bigotimes L_i$ sur $M_{0,I}$ : $3$ pour $L$ et $6$ pour $L_i$ quand $n = 4$ ; pour $n = 5$ et~$6$, les valeurs sont seulement posées en question. Mathlib n'a pas les espaces de modules $M_{0,n}$.\piste{146-equivariant-trivialising-exponents}
\end{itemize}
